\section{Chronology Is Part of the Evidence}

The Gospel chronology should not be recruited into a theory it does not state.
The Synoptics frame the preparations and meal in Passover language: the
disciples prepare the Passover, Jesus reclines with the Twelve, and the meal
opens directly into Gethsemane and the Passion (Matt 26:17--30; Mark
14:12--26; Luke 22:7--23). John's Gospel does not narrate the bread-and-cup
institution. It places the foot washing and farewell discourse before the
Passover feast and gives the Passion an intense paschal configuration (John
13:1; 18:28; 19:14, 31--36). The relation between the Synoptic and Johannine
calendars has generated several scholarly proposals. No proposed solution is
needed for the present doctrinal conclusion, and none warrants silently
rewriting one Evangelist through another.

All four narratives agree on the sequence that matters here: Jesus' final meal
with his disciples precedes arrest, trial, crucifixion, death, and burial. The
Resurrection follows. The Evangelists can present the meal as interpreting the
death and the death as fulfilling the meal without saying that the company at
table becomes chronologically simultaneous with the crucifixion. Narrative
sequence is not a crude obstacle that sacramental theology must overcome. It
is the created history in which the Word incarnate chose to institute,
suffer, die, and rise.

The institution narratives also differ in wording and arrangement. Matthew
and Mark emphasize covenant blood poured out for many and the eschatological
promise of the kingdom. Luke and Paul include the command of memorial and
explicitly connect the body with self-gift ``for you.'' Luke's longer cup
wording and Paul's received tradition closely align without becoming
identical transcripts. Catholic doctrine does not depend on collapsing them
into one reconstructed sentence. Their convergent grammar is sufficient:
Jesus takes, blesses or gives thanks, breaks, gives, identifies the gifts with
his Body and Blood, interprets his approaching death covenantally, and commands
an ecclesial repetition ordered to memory and expectation.

